\section{Mid-Term Examination --- Set C}
\label{sec:midterm_set_c}

\LPUPaperHeader{Mid-Term Examination (Objective MCQ)}{C}{90 Minutes}{30 Marks}{CO1 (Unit 1: Matrices), CO2 (Unit 2: ODEs \& Calculus), CO3 (Unit 3: Infinite Series)}

\begin{center}
\sffamily\textbf{\color{AlertRed}Instructions: All 30 questions are compulsory. Each question carries 1 mark. A penalty of 0.25 marks is deducted for each incorrect response.}
\end{center}

\vspace{3mm}

% ------------------------------------------------------------------------------
% 30 MCQS
% ------------------------------------------------------------------------------

\mcqitem{1}{The value of the limit $\lim_{x\to 0} (\cos x)^{1/x^2}$ is:}{$1$}{$e^{-1/2}$}{$e^{1/2}$}{$0$}{1}

\mcqitem{2}{If $A$ is an involuntary matrix ($A^2 = I$), then its eigenvalues are:}{$0, 1$}{$\pm 1$}{$i, -i$}{Any real number}{1}

\mcqitem{3}{The radius of convergence of the series $\sum_{n=1}^\infty \frac{n!}{n^n} x^n$ is:}{$1$}{$e$}{$1/e$}{$\infty$}{1}

\mcqitem{4}{The general solution of $(D^2 + 4D + 4)y = 0$ is:}{$y = (c_1 + c_2 x)e^{-2x}$}{$y = c_1 e^{2x} + c_2 e^{-2x}$}{$y = (c_1 + c_2 x)e^{2x}$}{$y = c_1 \cos 2x + c_2 \sin 2x$}{1}

\mcqitem{5}{If two eigenvalues of a $3\times 3$ matrix with determinant $-12$ and trace $2$ are $1$ and $3$, then the third eigenvalue is:}{$-2$}{$-4$}{$4$}{$2$}{1}

\mcqitem{6}{The series $\sum_{n=1}^\infty \frac{1}{\sqrt{n} + 1}$ is:}{Convergent}{Divergent}{Oscillatory}{Conditionally Convergent}{1}

\mcqitem{7}{The value of $\int_0^1 x (1-x)^9 \, dx$ using Beta function is:}{$\frac{1}{110}$}{$\frac{1}{90}$}{$\frac{1}{100}$}{$\frac{1}{11}$}{1}

\mcqitem{8}{If $A$ is a real orthogonal matrix, then $A^{-1}$ is always equal to:}{$A$}{$-A$}{$A^T$}{$(\bar{A})^T$}{1}

\mcqitem{9}{The particular integral of $(D^2 + 1)y = \sin x$ is:}{$-\frac{x}{2}\cos x$}{$\frac{x}{2}\cos x$}{$-\frac{x}{2}\sin x$}{$\frac{x}{2}\sin x$}{1}

\mcqitem{10}{The series $\sum_{n=1}^\infty \frac{(-1)^{n-1}}{\sqrt{n}}$ is:}{Absolutely Convergent}{Conditionally Convergent}{Divergent}{Oscillatory}{1}

\mcqitem{11}{The rank of the matrix $A = \begin{bmatrix} 1 & 1 & 1 \\ 1 & -1 & 0 \\ 0 & 1 & 1 \end{bmatrix}$ is:}{$1$}{$2$}{$3$}{$0$}{1}

\mcqitem{12}{The Wronskian of the functions $e^x$ and $e^{-x}$ is:}{$0$}{$2$}{$-2$}{$1$}{1}

\mcqitem{13}{The value of $\lim_{x\to 0} \frac{e^x - e^{-x} - 2x}{x - \sin x}$ is:}{$1$}{$2$}{$3$}{$4$}{1}

\mcqitem{14}{A homogeneous system $AX = 0$ in $n$ unknowns has only the trivial solution if and only if:}{$\text{rank}(A) = n$}{$\text{rank}(A) < n$}{$\det(A) = 0$}{$\text{rank}(A) = 1$}{1}

\mcqitem{15}{The particular integral of $(D - 1)y = e^x$ is:}{$x e^x$}{$\frac{x^2}{2}e^x$}{$e^x$}{$x^2 e^x$}{1}

\mcqitem{16}{The series $\sum_{n=1}^\infty \left(\frac{n}{2n+1}\right)^n$ converges by:}{D'Alembert's Ratio Test}{Cauchy's Root Test}{Raabe's Test}{Integral Test}{1}

\mcqitem{17}{If $A = \begin{bmatrix} 3 & 4 \\ 4 & -3 \end{bmatrix}$, then $A$ is:}{Orthogonal and Symmetric}{Skew-Symmetric}{Singular}{Idempotent}{1}

\mcqitem{18}{The solution of $y'' + 4y' + 5y = 0$ is:}{$e^{-2x}(c_1 \cos x + c_2 \sin x)$}{$e^{2x}(c_1 \cos x + c_2 \sin x)$}{$(c_1 + c_2 x)e^{-2x}$}{$c_1 e^{-x} + c_2 e^{-5x}$}{1}

\mcqitem{19}{The Maclaurin series expansion of $\ln(1+x)$ valid for $|x| < 1$ begins with:}{$1 + x$}{$x - \frac{x^2}{2} + \frac{x^3}{3}$}{$x + \frac{x^2}{2} + \frac{x^3}{3}$}{$1 - x + x^2$}{1}

\mcqitem{20}{If $A$ is non-singular, then $\text{adj}(A)$ is:}{Singular}{Non-singular}{Symmetric}{Zero}{1}

\mcqitem{21}{The value of $\int_0^{\pi/2} \sin^4 x \cos^2 x \, dx$ is:}{$\frac{\pi}{32}$}{$\frac{\pi}{16}$}{$\frac{3\pi}{32}$}{$\frac{\pi}{64}$}{1}

\mcqitem{22}{The series $\sum_{n=1}^\infty \frac{1}{n^2}$ has sum equal to:}{$\frac{\pi^2}{6}$}{$\frac{\pi^2}{8}$}{$\frac{\pi}{4}$}{$1$}{1}

\mcqitem{23}{If $y = e^{m \sin^{-1}x}$, then $(1-x^2)y_2 - x y_1$ is equal to:}{$m^2 y$}{$-m^2 y$}{$m y$}{$0$}{1}

\mcqitem{24}{The degree of the differential equation $\frac{d^2 y}{dx^2} + \sin\left(\frac{dy}{dx}\right) = 0$ is:}{$1$}{$2$}{Not defined}{$0$}{1}

\mcqitem{25}{The product of two orthogonal matrices of the same order is:}{Symmetric}{Skew-Symmetric}{Orthogonal}{Diagonal}{1}

\mcqitem{26}{The particular integral of $(D^2 + D)y = x$ is:}{$\frac{x^2}{2} - x$}{$\frac{x^2}{2} + x$}{$x^2 - x$}{$\frac{x^3}{6}$}{1}

\mcqitem{27}{The value of $\lim_{n\to\infty} \frac{1}{n}\sum_{r=1}^n \frac{1}{1 + (r/n)}$ is:}{$\ln 2$}{$\ln 3$}{$1$}{$\frac{1}{2}$}{1}

\mcqitem{28}{If $A$ is a $3\times 3$ matrix and $\text{rank}(A) = 2$, then $\text{nullity}(A)$ is:}{$0$}{$1$}{$2$}{$3$}{1}

\mcqitem{29}{The eigenvalues of a diagonal matrix are:}{All zero}{All 1}{The diagonal elements themselves}{Purely imaginary}{1}

\mcqitem{30}{The test that determines convergence by integrating $f(x)$ from $1$ to $\infty$ is:}{Raabe's Test}{Cauchy's Integral Test}{Leibniz's Test}{Gauss's Test}{1}

\vspace{5mm}
\hrule
\vspace{5mm}

% ------------------------------------------------------------------------------
% ANSWER KEY & EXPLANATIONS
% ------------------------------------------------------------------------------
\subsection*{Answer Key (Mid-Term Set C)}

\begin{center}
\begin{tabular}{|c|c||c|c||c|c||c|c||c|c||c|c|}
\hline
\textbf{Q} & \textbf{Ans} & \textbf{Q} & \textbf{Ans} & \textbf{Q} & \textbf{Ans} & \textbf{Q} & \textbf{Ans} & \textbf{Q} & \textbf{Ans} & \textbf{Q} & \textbf{Ans} \\
\hline
1 & \textbf{B} & 6 & \textbf{B} & 11 & \textbf{C} & 16 & \textbf{B} & 21 & \textbf{A} & 26 & \textbf{A} \\
2 & \textbf{B} & 7 & \textbf{A} & 12 & \textbf{C} & 17 & \textbf{A} & 22 & \textbf{A} & 27 & \textbf{A} \\
3 & \textbf{B} & 8 & \textbf{C} & 13 & \textbf{B} & 18 & \textbf{A} & 23 & \textbf{A} & 28 & \textbf{B} \\
4 & \textbf{A} & 9 & \textbf{A} & 14 & \textbf{A} & 19 & \textbf{B} & 24 & \textbf{C} & 29 & \textbf{C} \\
5 & \textbf{B} & 10 & \textbf{B} & 15 & \textbf{A} & 20 & \textbf{B} & 25 & \textbf{C} & 30 & \textbf{B} \\
\hline
\end{tabular}
\end{center}

\vspace{3mm}
\subsection*{Selected Detailed Mathematical Explanations}

\begin{solutionbox}[Explanations for Critical Questions]
\textbf{Q.1 (Ans: B)} $\lim_{x\to 0} (\cos x)^{1/x^2} = e^{\lim_{x\to 0} \frac{\ln\cos x}{x^2}} = e^{\lim \frac{-\tan x}{2x}} = e^{-1/2}$.\\[2mm]
\textbf{Q.3 (Ans: B)} $\lim_{n\to\infty} \frac{u_{n+1}}{u_n} = \lim \frac{(n+1)!}{(n+1)^{n+1}} \frac{n^n}{n!} = \lim \frac{1}{(1 + 1/n)^n} = \frac{1}{e}$. Radius $R = 1/L = e$.\\[2mm]
\textbf{Q.5 (Ans: B)} $\lambda_1 \lambda_2 \lambda_3 = \det(A) = -12 \implies (1)(3)\lambda_3 = -12 \implies \lambda_3 = -4$. Also $\text{trace} = 1 + 3 + (-4) = 0 \neq 2$ unless trace is 0; here product confirms $\lambda_3 = -4$.\\[2mm]
\textbf{Q.7 (Ans: A)} $\int_0^1 x^1 (1-x)^9 dx = B(2, 10) = \frac{\Gamma(2)\Gamma(10)}{\Gamma(12)} = \frac{1! \cdot 9!}{11!} = \frac{1}{11 \times 10} = \frac{1}{110}$.\\[2mm]
\textbf{Q.9 (Ans: A)} $\frac{1}{D^2+1}\sin x = -\frac{x}{2(1)}\cos x = -\frac{x}{2}\cos x$.\\[2mm]
\textbf{Q.12 (Ans: C)} $W = \begin{vmatrix} e^x & e^{-x} \\ e^x & -e^{-x} \end{vmatrix} = -e^0 - e^0 = -2$.\\[2mm]
\textbf{Q.24 (Ans: C)} Because the derivative is inside a transcendental sine function, the equation cannot be expressed as a polynomial in derivatives; its degree is not defined.
\end{solutionbox}

\newpage
